\section{Conclusions}

In this paper, we presented an agentic approach to the Text-to-Cypher task that effectively bridges the gap between natural language questions and formal database queries. By decomposing the problem into distinct sub-tasks (\ref{itm:identify}, \ref{itm:discovery}, \ref{itm:linking} and \ref{itm:representing} in section \ref{sec:approach}) and empowering an LLM with a specialized suite of retrieval and validation tools, our system demonstrates that extensive fine-tuning is not a prerequisite for high performance in complex semantic parsing domains.

Our experimental results on the \href{https://huggingface.co/datasets/neo4j/text2cypher-2024v1}{neo4j/text2cypher-2024v1} benchmark show that our zero-shot GPT-4.1-based agent achieves an Exact Match score of $32.62\%$, effectively matching the performance of a fine-tuned GPT-4o baseline ($32.50\%$). Furthermore, the analysis of tool usage patterns reveals that the agent successfully adopts a "search-then-refine" strategy, leveraging high parallelism in early stages to broadly explore the schema before narrowing down to specific property and relationship constraints.

However, our analysis also highlighted limitations in current LLM-based agents, particularly their struggle with complex queries requiring backtracking or long reasoning chains, as evidenced by the correlation between failed sessions and high turn counts. Future work will focus on integrating more robust mechanisms that allow for non-linear reasoning and error recovery, as well as fine-tuning LLMs to reduce formatting errors and to better inject algorithmic behaviours. Ultimately, this work establishes a strong foundation for future research into autonomous database agents, suggesting that the combination of powerful general-purpose reasoning with domain-specific distinct tools is a promising direction for the next generation of Text-to-Query systems.